\makeatletter
\def\rap@fontpath{../../../Common/}
\makeatother
\documentclass{../../../Common/Latex/Templates/rapport}
% ---------------------------------------------------------------------------
% Shared settings for the H2026 General Assembly document set.
%
% Usage, directly after \documentclass{../../../Common/Latex/Templates/rapport}
% in a document one folder below this file:
%   % ---------------------------------------------------------------------------
% Shared settings for the H2026 General Assembly document set.
%
% Usage, directly after \documentclass{../../../Common/Latex/Templates/rapport}
% in a document one folder below this file:
%   % ---------------------------------------------------------------------------
% Shared settings for the H2026 General Assembly document set.
%
% Usage, directly after \documentclass{../../../Common/Latex/Templates/rapport}
% in a document one folder below this file:
%   \input{../ga-common}
%
% Image paths are relative to the compiled document's folder, so every
% document using this file must sit one folder below it.
% ---------------------------------------------------------------------------

% Meeting details, used on covers and in each document's ID block.
\newcommand{\meetingdate}{30 September 2026}
\newcommand{\meetingtime}{12:00}
\newcommand{\meetinglocation}{Festsalen, NMBU}
\newcommand{\meetingline}{\textbf{Meeting:} \meetingdate, \meetingtime, \meetinglocation}

% Header: moon logo on the left, full-text logo on the right. \projectname is
% still set because the cover title uses it.
\projectname{Eclipse NMBU}
\projectlogo{../../../Common/Eclipse Images/EclipseMoon}
\orglogo{../../../Common/Eclipse Images/EclipseFullText}

%
% Image paths are relative to the compiled document's folder, so every
% document using this file must sit one folder below it.
% ---------------------------------------------------------------------------

% Meeting details, used on covers and in each document's ID block.
\newcommand{\meetingdate}{30 September 2026}
\newcommand{\meetingtime}{12:00}
\newcommand{\meetinglocation}{Festsalen, NMBU}
\newcommand{\meetingline}{\textbf{Meeting:} \meetingdate, \meetingtime, \meetinglocation}

% Header: moon logo on the left, full-text logo on the right. \projectname is
% still set because the cover title uses it.
\projectname{Eclipse NMBU}
\projectlogo{../../../Common/Eclipse Images/EclipseMoon}
\orglogo{../../../Common/Eclipse Images/EclipseFullText}

%
% Image paths are relative to the compiled document's folder, so every
% document using this file must sit one folder below it.
% ---------------------------------------------------------------------------

% Meeting details, used on covers and in each document's ID block.
\newcommand{\meetingdate}{30 September 2026}
\newcommand{\meetingtime}{12:00}
\newcommand{\meetinglocation}{Festsalen, NMBU}
\newcommand{\meetingline}{\textbf{Meeting:} \meetingdate, \meetingtime, \meetinglocation}

% Header: moon logo on the left, full-text logo on the right. \projectname is
% still set because the cover title uses it.
\projectname{Eclipse NMBU}
\projectlogo{../../../Common/Eclipse Images/EclipseMoon}
\orglogo{../../../Common/Eclipse Images/EclipseFullText}

\usepackage{array}
\usepackage{tabularx}

\reporttitle{Case 7: Election of Board Members}

\begin{document}

\reportmainmatter
\small
\setlength{\parindent}{0pt}
\setlength{\tabcolsep}{6pt}
\renewcommand{\arraystretch}{1.2}

\begin{center}
  {\Large\bfseries Election of Board Members (valg av styremedlemmer): CEO, CTO, CFO, CMO}\par
  \vspace{0.15cm}
  {\normalsize 26H-ECL-GA-CASE6B \textperiodcentered{} Case 7 \textperiodcentered{} Proposed by: the Board}\par
  \vspace{0.1cm}
  {\small \meetingline}\par
\end{center}

\vspace{0.3cm}
\begin{tabularx}{\textwidth}{|>{\bfseries\raggedright\arraybackslash}p{4.2cm}|X|}
  \hline
  C-suite deliberation & The C-suite has reviewed the proposal and recommends it to the General Assembly for approval.                                                                                                                                                                                                       \\
  \hline
  Background           & Old Bylaws \S4.7: all Board members are elected at the General Assembly, by majority. The Bylaw Amendment was not adopted under Case 5, and Case 6 (26H-ECL-GA-CASE5B) fixed the four roles as CEO, CTO, CFO and CMO, as defined in Old Bylaws \S5.1--\S5.5.                                        \\
  \hline
  Proposal             & The General Assembly resolves to elect, one role at a time in the order CEO, CTO, CFO, CMO, the members recorded as elected in the Meeting Protocol (26H-ECL-GA-PROT1), for a term of two student semesters (Old Bylaws \S5.2--\S5.5), using the election mechanics in Old Bylaws \S4.7.1--\S4.7.2. \\
  \hline
  External information & Procedural Case: no Board recommendation on individual candidates. Attachments: C-suite structure (26H-ECL-GA-CASE5B), Candidate Documentation (26H-ECL-GA-CAND1), Election Ballot (26H-ECL-GA-BALLOT2B).                                                                                           \\
  \hline
\end{tabularx}

\end{document}
